\section{From Section 3.6 to the integrated experiments}
\label{sec:history}

The earliest question was whether an inexpensive smooth observation
functional could provide useful score information alongside the GenUT
filter. The Section 3.6 sidecar was designed as a separate deterministic
proposal force, with the canonical value retained for endpoint correction.
That was a legitimate mechanics experiment, but it was not a construction of
the canonical total score and it never supplied a known-mean control variate.

Its first repair concerned observation timing. Let $w_i^-$ be normalized
weights before the current observation, $g_i$ its likelihood at particle $i$,
and $Z=\sum_i w_i^-g_i$. The posterior weights are $w_i^+=w_i^-g_i/Z$.
The historical capture hook ran before reset but after this observation
update. Reusing its $w_i^+$ in a second likelihood average therefore gave
\begin{equation}
 \widetilde Z=\sum_iw_i^+g_i
 =\frac{\sum_iw_i^-g_i^2}{Z},\qquad
 \widetilde Z-Z=\frac{\operatorname{Var}_{w^-}(g)}{Z}.
 \label{eq:historical-double-observation}
\end{equation}
For two equally weighted particles with likelihoods 1 and 3, the intended
increment is 2, while the reused posterior weights give 2.5. The word
``pre-reset'' did not make the captured weights pre-observation. The repaired
sidecar used the actual prior observation weights and differentiated its
declared observation covariance dependence, while holding the captured cloud,
weights, and kernel fixed. It remained a conditional derivative.

The September 3 pilot used one-step diagonal LGSSMs: $(d,N,T)=(1,4,1)$ and
$(3,6,1)$. Four calibration paths and eight streams per scope selected the
grid boundary $\rho=0.8$; eight disjoint validation paths and eight streams
per scope supplied the holdout comparison. The recorded relative
path-level squared-bias changes were 6.6082 with interval
$[-0.2825,19.4854]$ and 0.2209 with interval $[-0.9342,1.5096]$.
Both intervals crossed zero, and the recorded stream-MCSE ratios, ranging
from 0.1756 to 0.4709, failed the threshold 0.10. These tiny CPU/non-XLA
diagnostics supplied no ranking or bias-reduction evidence. Separate
reversibility and volume checks passed at about machine precision, but
proposal mechanics did not establish score accuracy, HMC mixing, or posterior
correctness. The sidecar remained isolated and unpromoted.

The next stages sought to remove the conditional-derivative restriction.
Phases 1--3 implemented complete mixture differentials, a fixed-chart support
reference, and an auxiliary functional constructed from the actual canonical
trace. Phase 4A then fed kernelized observation weights and their tangents
through the complete reset. Phase 4B added the source paper's mixture
resampling mechanism after Contract--E and carried the resulting weights and
covariance marks into later time steps. These changes made the full feedback
testable. They also defined new finite programs; no residual correction made
them equal to the original atom program.

The chronology therefore does not support a claim that a control variate was
completed and then simplified. The executed work moved from an auxiliary
conditional functional to separately named integrated experiments. This was
useful for determining whether those mechanisms improved model-score error.
It did not solve the original target-preserving control-variate problem.
The negative results below are the reason to reconsider the insertion point
and the score estimator together.
